\section{Introduction}
\label{sec:introduction}

Tropical peat swamp forests are among the most carbon-dense terrestrial ecosystems on the planet, sequestering up to twenty times more organic carbon per unit area than adjacent upland tropical rainforests on mineral soils \citep{Page2002Nature, Dommain2011ProcRSocB}. In undisturbed equatorial environments, high phreatic water tables remain near or above the surface throughout the annual cycle, creating anoxic conditions that retard organic matter decomposition and render natural peat formations virtually non-combustible \citep{Wosten2008Geoderma, Dadap2021NatComm}. Over the past three decades, however, Southeast Asian peatlands across Sumatra and Borneo have experienced severe anthropogenic disturbance from industrial logging, agricultural conversion, and massive state-sponsored drainage schemes \citep{Miettinen2016ERL, Dohong2017RestorEcol, Langner2007IJWF, Carlson2012PNAS}.

In Central Kalimantan, Indonesia, the Mega Rice Project (ex-PLG) excavated over 4,000\,km of primary, secondary, and tertiary drainage canals between 1995 and 1999 across more than one million hectares of ombrotrophic peat domes between the Sebangau, Kahayan, Kapuas, and Barito rivers \citep{Silvius2005Wetlands, Ritzema2014HESS}. These canals accelerate gravitational drainage during dry seasons, drawing groundwater tables below the regulatory limit of $-40$\,cm established under Indonesian Government Regulations (PP No.~71/2014 and PP No.~57/2016). When dry spells coincide with positive El Ni\~no--Southern Oscillation (ENSO) anomalies or positive Indian Ocean Dipole (pIOD) events, drained peat ignites readily \citep{Cochrane2003Nature, Hooijer2012Biogeosciences, Field2016PNAS}. Once ignited, smoldering combustion propagates horizontally and vertically through the catotelm for months, releasing ancient carbon sequestered over millennia \citep{Usup2004Tropics, Wooster2018NatGeosci, Vetrita2020RemoteSensing}.

Major fire episodes in 1997, 2015, and 2019 released an estimated 0.81--2.57\,Gt, 1.2\,Gt, and 0.71\,Gt of $\text{CO}_2$-equivalent emissions, respectively, at times exceeding the daily fossil fuel emissions of major industrialized nations \citep{Page2002Nature, Field2016PNAS, VanderWerf2017ESSD}. Peat smoldering also generates fine particulate matter ($\text{PM}_{2.5}$), carbon monoxide, and aerosol precursors that form regional smoke-haze across Indonesia, Singapore, and Malaysia, causing acute public health emergencies and substantial economic losses \citep{Gaveau2014PLOSONE, Kiely2021ACP, Graham2022SciRep}.

In the domain of land use policy, environmental governance, and applied spatial geography, resolving tropical peatland wildfire crises encounters two critical impasses:

First, \textbf{the territorial political economy of land tenure culpability and concession frontier spillover remains clouded by empirical ambiguity and legal asymmetry}. Public discourse and regulatory enforcement in Indonesia have centered primarily on corporate concessionaires, anchored in the strict statutory liability framework of Law No.~32/2009 on Environmental Protection and Management (Article 88), which imposes no-fault liability for any fire occurring within certified concession perimeters \citep{Tacconi2003CIFOR, Carlson2012PNAS}. However, historical state territorialization partitioned Indonesia's land estate into the State Forest Estate (\textit{Kawasan Hutan}) and Other Utilization Areas (\textit{Area Penggunaan Lain}, APL). In post-Mega Rice Project Central Kalimantan, vast expanses of drained peatlands were excised from the forest estate into APL, but left unallocated without designated concessionaires, Forest Management Units (KPH), or dedicated fire suppression mandates \citep{Cattau2016GCB, Putranto2023LandUsePol}. Concurrently, artificial drainage networks constructed within corporate concessions inevitably lower regional groundwater tables across contiguous hydrological domes, desorbing adjacent state peatlands beyond concession borders. Whether corporate strict liability deters fire inside borders while inadvertently displacing ignitions onto adjacent unmanaged APL lands---creating a ``frontier governance void''---remains untested through rigorous econometric or spatial boundary discontinuity techniques.

Second, \textbf{unmitigated spatial autocorrelation induces insidious spatial memorization in predictive spatial land use models}. Machine learning classifiers, including Random Forests, are increasingly deployed for spatial wildfire susceptibility and active combustion mapping \citep{Breiman2001ML}. However, terrestrial environmental processes strictly operate under Tobler's First Law of Geography: everything is related to everything else, but near things are more related than distant things \citep{Tobler1970EconGeogr, Legendre1993Ecology}. When spatial hazard models ingest geodetic coordinates (\texttt{latitude} and \texttt{longitude}) or administrative unit labels as independent predictors, recursive tree partitioning algorithms exploit spatial coordinates as empirical shortcuts, memorizing historical spatial clusters rather than learning underlying biophysical combustion mechanics \citep{Meyer2018EcolModel, Dormann2007Ecography}. Under standard random cross-validation, spatial autocorrelation between geographically proximate training and validation samples causes severe data leakage \citep{Roberts2017Ecography, Wadoux2021EcolModel}. Consequently, while coordinate-dependent models report deceptively high predictive skill within sampled neighborhoods, their transferability collapses precipitously when deployed across out-of-sample administrative jurisdictions or unfamiliar peat dome systems \citep{MeyerPebesma2021NatComm, Schratz2019EcolModel}.

To resolve these fundamental land tenure and spatial methodological dilemmas, this study presents an empirical investigation of the wildfire crisis in Central Kalimantan. Utilizing 43,937 satellite active fire observations from NASA FIRMS during the acute October 2026 fire pulse, we test two central land governance hypotheses:
\begin{enumerate}
    \item[$\mathbf{H_1}$:] \textit{Corporate strict liability creates a sharp spatial discontinuity at concession perimeters: fires are actively suppressed within certified concession boundaries, but active fire density escalates sharply directly across the cadastral frontier onto unallocated APL state land due to drainage spillover and an open-access institutional governance void.}
    \item[$\mathbf{H_2}$:] \textit{A pure biophysical feature architecture free from spatial coordinates generalizes across distinct geographical blocks by learning transferable environmental, hydrological, and edaphic combustion thresholds, whereas coordinate-dependent models suffer spatial memorization collapse under Spatial Block Cross-Validation.}
\end{enumerate}

To evaluate these hypotheses within a land use policy framework, we: (1) analyze decadal burned area records (2018--2026) from the Ministry of Environment and Forestry (SiPongi KLHK) alongside micro-scale diurnal satellite dynamics; (2) map 13 primary canal vectors in UTM Zone 49S projection across the ex-Mega Rice Project landscape; (3) perform a comprehensive spatial overlay against 812 cadastral concessions (627 oil palm HGU and 185 commercial forestry PBPH) to resolve land tenure attribution; (4) deploy a Spatial Boundary Discontinuity Design (Spatial BDD) to quantify the boundary step jump and distance-decay spillover within a 3-km frontier corridor; (5) benchmark coordinate-dependent (Model~A) and pure biophysical (Model~B) Random Forest models using 5-fold Spatial Block Cross-Validation (\texttt{GroupKFold} by regency); and (6) formulate actionable land governance reforms to close the concession frontier governance void.
